\section{The bipolar model as coded}
\label{sec:model}

\paragraph*{States, regimes, and strategies.}
Two countries, $i$ and $j$, interact in discrete time and discount the future
at rate $\delta$. The state $s \in \{\hI, \dI, \dJ, \hJ\}$ records who leads:
$i$ is the hegemon at $\hI$, $i$ is the dominant country of a bipolar world
at $\dI$, and symmetrically for $j$. The code indexes the states $1$ to $4$
in this order. The regime $\rho \in \{P, R\}$ (power or rules) is a second
state variable that persists until a country changes it. Under $P$ at $\hI$,
the hegemon builds rules with probability $\lambda_1$ and pays the building
cost $c$ once. Under $R$ at $\dI$, the dominant country keeps the rules with
probability $\lambda_2$; breaking them costs nothing and returns the regime
to $P$. Country $j$ uses the same $(\lambda_1, \lambda_2)$ at $\hJ$ and
$\dJ$, so a symmetric Markov equilibrium is a pair $(\lambda_1, \lambda_2)$.
Two further restrictions are imposed as constraints: a hegemon never
abandons rules that already exist, and no country builds rules in a
dominance state.

\paragraph*{Payoffs.}
Country $i$'s per-period payoff depends on the state and the regime.
Under rules it equals $u_R = 1$ in every state. Under power it equals
\begin{align*}
  u^P(\hI) = 0, \qquad u^P(\dI) = 1 + k, \qquad u^P(\dJ) = 1 - 2k, \qquad u^P(\hJ) = 0,
\end{align*}
with $k = 0.1$ in every figure. The joint per-period payoff is therefore
$2$ under rules, $2 - k = 1.9$ in a dominance state under power (the
collective-action loss), and $0$ in a hegemonic state under power. The
building cost $c$ is the horizontal axis of most maps.

\paragraph*{Transitions.}
The state follows a Markov chain with transition matrix $Q$, rows and
columns ordered $(\hI, \dI, \dJ, \hJ)$:
\begin{align*}
  Q = \begin{pmatrix}
    0.10 & 0.90 & 0    & 0    \\
    0.45 & 0.10 & 0.45 & 0    \\
    0    & 0.45 & 0.10 & 0.45 \\
    0    & 0    & 0.90 & 0.10
  \end{pmatrix}.
\end{align*}
Hegemony is followed by the dominance of the same country with probability
$0.9$, and a dominant country either becomes hegemon or loses its lead with
equal probability $0.45$. The persistence of the dominance states,
$q_{\dI,\dI} = q_{\dJ,\dJ} = 0.1$, is the parameter varied in Figure 1.B.3.
Transitions do not depend on the regime except in Figure 1.F.1, where
choosing power in a dominance state raises the chance of staying dominant.

\paragraph*{Value functions.}
Let $V^{\rho}_s$ denote $i$'s value at state $s$ under regime $\rho$ at the
start of a period, $V^{\rho} = (V^{\rho}_1, \dots, V^{\rho}_4)'$, $Q_s$ the
$s$-th row of $Q$, and $Q^P_s$ the corresponding row under power ($Q^P = Q$
except in Figure 1.F.1). The scripts solve the system
\begin{align*}
  V^P_1 &= \lambda_1 \big[u^R_1 - c + \delta Q_1 V^R\big] + (1 - \lambda_1) \big[u^P_1 + \delta Q^P_1 V^P\big], \\
  V^R_1 &= u^R_1 + \delta Q_1 V^R, \\
  V^P_2 &= u^P_2 + \delta Q^P_2 V^P, \\
  V^R_2 &= \lambda_2 \big[u^R_2 + \delta Q_2 V^R\big] + (1 - \lambda_2) \big[u^P_2 + \delta Q^P_2 V^P\big], \\
  V^P_3 &= u^P_3 + \delta Q^P_3 V^P, \\
  V^R_3 &= \lambda_2 \big[u^R_3 + \delta Q_3 V^R\big] + (1 - \lambda_2) \big[u^P_3 + \delta Q^P_3 V^P\big], \\
  V^P_4 &= \lambda_1 \big[u^R_4 + \delta Q_4 V^R\big] + (1 - \lambda_1) \big[u^P_4 + \delta Q^P_4 V^P\big], \\
  V^R_4 &= u^R_4 + \delta Q_4 V^R,
\end{align*}
where $u^{\rho}_s$ is $i$'s payoff at state $s$ under $\rho$. At $\hJ$ it is
$j$ who pays the building cost, so $i$'s value under power at $\hJ$ carries
no cost term.

\paragraph*{Equilibrium conditions.}
Write the one-shot deviation values at $\hI$ under $P$ (build or not) and at
$\dI$ under $R$ (keep or break) as
\begin{align*}
  D^P_1 &= u^P_1 + \delta Q^P_1 V^P, & D^R_1 &= u^R_1 - c + \delta Q_1 V^R, \\
  D^P_2 &= u^P_2 + \delta Q^P_2 V^P, & D^R_2 &= u^R_2 + \delta Q_2 V^R.
\end{align*}
A pair $(\lambda_1, \lambda_2)$ together with $(V^P, V^R)$ is an equilibrium
when it solves the system above and satisfies four conditions:
\begin{enumerate}
  \item[(a)] $\lambda_1 = 1$ if $D^R_1 > D^P_1$, $\lambda_1 = 0$ if $D^R_1 < D^P_1$, and $\lambda_1 \in (0,1)$ only if $D^R_1 = D^P_1$;
  \item[(b)] the same for $\lambda_2$ with $D^R_2$ and $D^P_2$;
  \item[(c)] $D^P_1 \leq u^R_1 + \delta Q_1 V^R$: at $\hI$ under $R$ the hegemon keeps the rules;
  \item[(d)] $D^P_2 \geq u^R_2 - c + \delta Q_2 V^R$: at $\dI$ under $P$ the dominant country does not build rules.
\end{enumerate}
Conditions (a) and (b) are the optimality of the two probabilities; (c) and
(d) are the two restrictions on off-path behavior.

\paragraph*{Classification of equilibria.}
The maps color each grid point by the equilibrium pair:
$(\lambda_1, \lambda_2) = (1, 1)$ is a rules-based equilibrium (the hegemon
builds rules and the dominant country keeps them); $(1, 0)$ is a cyclic
equilibrium (rules are built under hegemony and broken as soon as the world
turns bipolar); $(0, 1)$ is a rules-compatible power-based equilibrium
(nobody builds, but rules would be kept if they existed); $(0, 0)$ is a
rules-resistant power-based equilibrium (nobody builds and rules would be
broken). Three-class maps merge the last two classes into ``power-based
equilibria''.

\paragraph*{Two-leader variant (Figure 6 of the main paper).}
The script behind Figure 6 of the main paper adds a fifth state $N$ between $\dI$ and
$\dJ$: a multipolar state in which neither leader dominates and both act at
once. The transition matrix, with rows and columns ordered $(\hI, \dI, N,
\dJ, \hJ)$, keeps the baseline rows for the hegemonic states and gives a
dominant country equal chances $0.45$ of becoming hegemon and of falling
into $N$, where the world stays with probability $0.1$ or moves to either
dominance state with probability $0.45$:
\begin{align*}
  Q = \begin{pmatrix}
    0.10 & 0.90 & 0    & 0    & 0    \\
    0.45 & 0.10 & 0.45 & 0    & 0    \\
    0    & 0.45 & 0.10 & 0.45 & 0    \\
    0    & 0    & 0.45 & 0.10 & 0.45 \\
    0    & 0    & 0    & 0.90 & 0.10
  \end{pmatrix}.
\end{align*}
At $N$ the payoff depends on the pair of actions: $1$ when both play rules,
$1 - 2k = 0.8$ when $i$ plays rules and $j$ power, $1 + k = 1.1$ in the
opposite case, and $0.95$, the average of the two asymmetric outcomes, when
both play power; under regime $P$ both play power. A third probability
$\lambda_N$ gives the chance that a leader plays rules at $N$ under regime
$R$, with the joint-support assumption that rules survive there only if both
leaders play rules, so the continuation under rules carries probability
$\lambda_N^2$. The value equations and the four conditions above extend
accordingly, with two additions: $\lambda_N$ is optimal given the other
leader's $\lambda_N$, and nobody builds rules at $N$ under $P$. The script
keeps, among multiple solutions, the one with $\lambda_2 = \lambda_N = 1$.
In the three-class map, rules-based equilibria are $(\lambda_1, \lambda_2,
\lambda_N) = (1, 1, 1)$, cyclic equilibria have $\lambda_1 = 1$ and rules
broken at $\dI$ or at $N$, and power-based equilibria have $\lambda_1 = 0$.
